\documentclass[a4,11pt]{aleph-notas}

% -- Paquetes adicionales
\usepackage{enumitem}
\usepackage{aleph-comandos}

% -- Datos
\institucion{Facultad de Ciencias Exactas, Naturales y Ambientales}
\carrera{Catálogo STEM}
\asignatura{Matemática III}
\tema{Resumen 03: Topología de $\mathbb{R}^n$}
\autor[A. Merino]{Andrés Merino}
\fecha{Periodo 2026-1}

\logouno[0.16\textwidth]{Logos/logoPUCE_04_ac}
\definecolor{colortext}{HTML}{1957d1}
\definecolor{colordef}{HTML}{1957d1}
\fuente{montserrat}


\begin{document}

\encabezado

En este resumen consideraremos $n>1$.

%%%%%%%%%%%%%%%%%%%%%%%%%%%%%%%%%%%%%%
\section{Bolas y conjuntos abiertos y cerrados}
%%%%%%%%%%%%%%%%%%%%%%%%%%%%%%%%%%%%%%

\begin{defi}[Bola abierta]
    Dados $x_0\in\R^n$ y $r>0$, se define la bola abierta de centro $x_0$ y radio $r$ como el conjunto
    \[
        B(x_0,r)=\{x\in\R^n : \|x-x_0\|<r\}.
    \]
\end{defi}

\begin{defi}[Conjunto abierto y cerrado]
    Dado un conjunto $A\subset \R^n$, se dice que es \textbf{abierto} si para todo $x\in A$, existe $r>0$ tal que
    \[
        B(x,r)\subset A.
    \]
    Por otro lado, se dice que $A$ es \textbf{cerrado} si $A^c$ es abierto.
\end{defi}

%%%%%%%%%%%%%%%%%%%%%%%%%%%%%%%%%%%%%%
\section{Clausura, frontera e interior}
%%%%%%%%%%%%%%%%%%%%%%%%%%%%%%%%%%%%%%

\begin{defi}[Clausura, derivado, frontera e interior]
    Dado un conjunto $A\subset \R^n$, se define:
    \begin{itemize}
        \item la \textbf{clausura} de $A$:
        \[
            \cl{A}=\{x\in \R^n: (\forall r>0)(B(x,r)\cap A\neq \emptyset)\};
        \]
        \item el \textbf{derivado} de $A$:
        \[
            A'=\{x\in \R^n: (\forall r>0)((B(x,r)\setminus\{x\})\cap A\neq \emptyset)\};
        \]
        \item la \textbf{frontera} de $A$:
        \[
            \delta(A)=\{x\in \R^n: (\forall r>0)(B(x,r)\cap A\neq \emptyset\ \land\ B(x,r)\cap A^c\neq \emptyset)\};
        \]
        \item el \textbf{interior} de $A$:
        \[
            \inte(A)=\{x\in \R^n: (\exists r>0)(B(x,r)\subset A)\}.
        \]
    \end{itemize}
\end{defi}

\newpage
%%%%%%%%%%%%%%%%%%%%%%%%%%%%%%%%%%%%%%
\section{Conjunto acotado}
%%%%%%%%%%%%%%%%%%%%%%%%%%%%%%%%%%%%%%

\begin{defi}[Conjunto acotado]
    Dado un conjunto $A\subset \R^n$, se dice que $A$ es \textbf{acotado} si existe $R>0$ tal que
    \[
        \|x\|<R,
    \]
    para todo $x\in A$.
\end{defi}

\end{document}
